% TODO checklist:
% - Inspected: src/config.py (APP_ENV handling, security settings)
% - Inspected: docker-compose.override.yml (development overrides)
% - Inspected: docker-compose.override.dev.yml (development settings)
% - Inspected: Q&A.md sections on deployment assumptions
% - Inspected: src/logging_setup.py (environment-specific logging)
% - Inspected: Makefile (up-cpu, up-gpu targets)

\section{Environment Profiles (Development, Testing, Production)}
\label{sec:env-profiles}

The Cineca Agentic Platform supports multiple environment profiles that adapt configuration, security settings, and operational behavior to the deployment context. This section describes the three primary environment profiles---development, testing, and production---and explains how the platform differentiates behavior across these environments.

\subsection{Environment Mode Detection}

The platform determines its operating environment through the \texttt{APP\_ENV} configuration variable, which accepts the following values:

\begin{itemize}
    \item \texttt{dev} (default): Development mode with relaxed security, full documentation, and debug-friendly logging
    \item \texttt{test}: Testing mode with mocked providers, reduced timeouts, and deterministic behavior
    \item \texttt{stage}/\texttt{staging}: Staging environment with production-like settings but relaxed security
    \item \texttt{prod}/\texttt{production}: Production mode with full security hardening and restricted access
\end{itemize}

The configuration system normalizes environment values to ensure consistent handling:

\begin{lstlisting}[language=Python, caption=Environment normalization in src/config.py]
@field_validator("APP_ENV", mode="before")
@classmethod
def normalize_env(cls, v: str) -> str:
    return (v or "dev").lower()
\end{lstlisting}

\subsection{Development Environment}

The development environment prioritizes developer experience, debugging capability, and rapid iteration. Key characteristics include:

\paragraph{Full API Documentation}
The Swagger UI and ReDoc documentation interfaces are enabled by default:

\begin{lstlisting}[language=Python, caption=Documentation enablement in development]
ENABLE_DOCS: bool = Field(
    default=True, 
    description="Enable Swagger/Redoc in non-prod"
)
\end{lstlisting}

\paragraph{Console Logging}
Development environments use human-readable console logging rather than JSON-structured logs:

\begin{lstlisting}[language=Python, caption=Environment-specific logging format]
def _wants_json() -> bool:
    """Return True if JSON logging should be used."""
    env = os.getenv("APP_ENV", "dev").strip().lower()
    return env in {"prod", "production"}  # JSON logs in production only
\end{lstlisting}

\paragraph{Relaxed Security Settings}
Development mode uses relaxed security defaults to facilitate local testing:

\begin{lstlisting}[language=bash, caption=Development environment overrides from docker-compose.override.yml]
services:
  app:
    command: uvicorn src.app:app --reload --host 0.0.0.0 --port 8000
    environment:
      PYTHONUNBUFFERED: '1'
      ENABLE_DOCS: '1'
      APP_ENV: 'dev'
      DEMO_MODE: 'true'
      RATE_LIMIT_MODE: 'test'  # Relaxed rate limits
      INTERNAL_TOKEN_MAX_TTL_SECONDS: '86400'  # Allow 24h tokens
\end{lstlisting}

\paragraph{Hot Reload}
The uvicorn server runs with \texttt{--reload} enabled, automatically restarting when source files change.

\paragraph{Demo Mode}
When \texttt{DEMO\_MODE=true}, the platform provides fallback responses when no LLM providers are registered, enabling UI development without requiring a running LLM.

\subsection{Testing Environment}

The testing environment is optimized for automated testing with predictable behavior and reduced resource requirements:

\paragraph{Reduced Timeouts}
The compute configuration module provides test-specific timeouts:

\begin{lstlisting}[language=Python, caption=Test mode timeout configuration]
@property
def recommended_step_timeout(self) -> int:
    if self.memgraph_nl_test_mode:
        return 90   # 90s for simple NL-to-Cypher queries
    if self.test_mode:
        return 60   # 60s for general test mode
    # ... device-specific defaults
\end{lstlisting}

\paragraph{Mocked Providers}
Tests can enable \texttt{SEED\_DEMO\_PROVIDER=false} to prevent automatic provider creation, allowing tests to configure their own mock providers.

\paragraph{Deterministic Memgraph Responses}
The \texttt{MEMGRAPH\_NL\_TEST\_MODE} flag enables deterministic testing of the NL-to-Cypher pipeline by using a JSON mapping file rather than LLM generation:

\begin{lstlisting}[language=bash, caption=Test mode configuration for Memgraph NL]
# Enable deterministic NL-to-Cypher testing
MEMGRAPH_NL_TEST_MODE=true
MEMGRAPH_RESPONSE_MODE=fallback-only
\end{lstlisting}

\paragraph{In-Memory Job Store}
Tests default to the in-memory job store to avoid Redis dependencies:

\begin{lstlisting}[language=Python, caption=Job store backend configuration]
JOB_STORE_BACKEND: str = Field(
    default="memory",
    description="'memory' or 'redis'. Use 'redis' for production."
)
\end{lstlisting}

Table~\ref{tab:test-env-settings} summarizes key testing environment settings.

\begin{table}[ht]
\centering
\caption{Testing environment configuration settings.}
\label{tab:test-env-settings}
\begin{tabular}{llp{6cm}}
\hline
\textbf{Setting} & \textbf{Value} & \textbf{Purpose} \\
\hline
LLM\_TEST\_MODE & true & Enables reduced timeouts and lightweight models \\
MEMGRAPH\_NL\_TEST\_MODE & true & Uses deterministic NL-to-Cypher mappings \\
MEMGRAPH\_RESPONSE\_MODE & fallback-only & Skips LLM summarization for speed \\
JOB\_STORE\_BACKEND & memory & Avoids Redis dependency \\
RATE\_LIMIT\_MODE & test & Relaxes rate limiting for test concurrency \\
DEMO\_MODE & true & Provides fallback responses without LLM \\
\hline
\end{tabular}
\end{table}

\subsection{Production Environment}

The production environment enforces strict security measures, optimizes for reliability, and enables comprehensive observability:

\paragraph{Security Headers}
Production mode enables security headers middleware:

\begin{lstlisting}[language=Python, caption=Production security header settings]
ENABLE_SECURITY_HEADERS: bool = Field(
    default=True, 
    description="Enable security headers middleware"
)
ENABLE_HSTS: bool = Field(
    default=True, 
    description="Enable Strict-Transport-Security header"
)
HSTS_MAX_AGE: int = Field(
    default=31536000,  # 1 year
    description="HSTS max-age in seconds"
)
\end{lstlisting}

\paragraph{Secure Cookie Configuration}
Cookies are configured with security attributes:

\begin{lstlisting}[language=Python, caption=Production cookie settings]
SECURE_COOKIES: bool = Field(
    default=False,  # Set True in production!
    description="Enable secure flag on cookies"
)
TRUST_PROXY: bool = Field(
    default=False,  # Set True behind reverse proxy
    description="Trust X-Forwarded-* headers"
)
\end{lstlisting}

\paragraph{JSON-Structured Logging}
Production environments use JSON-formatted logs for integration with log aggregation systems:

\begin{lstlisting}[language=Python, caption=Production logging configuration]
# From src/logging_setup.py
if _wants_json():
    # Configure structlog for JSON output
    processors.append(structlog.processors.JSONRenderer())
else:
    # Console-friendly output for development
    processors.append(structlog.dev.ConsoleRenderer())
\end{lstlisting}

\paragraph{Restricted API Documentation}
Production deployments can restrict access to API documentation:

\begin{lstlisting}[language=Python, caption=Documentation access control]
DOCS_AUTH: str = Field(
    default="internal",
    description="internal|public - controls Swagger 'Try it out'"
)
\end{lstlisting}

\paragraph{Redis Job Store}
Production deployments use Redis for persistent job storage:

\begin{lstlisting}[language=bash, caption=Production job store configuration]
JOB_STORE_BACKEND=redis
USE_POSTGRES_JOBS=true
\end{lstlisting}

\subsection{Environment Comparison Matrix}

Table~\ref{tab:env-comparison} provides a comprehensive comparison of settings across environment profiles.

\begin{table}[ht]
\centering
\caption{Configuration differences across environment profiles.}
\label{tab:env-comparison}
\small
\begin{tabular}{llll}
\hline
\textbf{Setting} & \textbf{Development} & \textbf{Testing} & \textbf{Production} \\
\hline
APP\_ENV & dev & test & prod \\
ENABLE\_DOCS & true & true & false/restricted \\
LOG\_LEVEL & DEBUG & INFO & WARNING \\
Log Format & Console & Console & JSON \\
ENABLE\_SECURITY\_HEADERS & false & false & true \\
ENABLE\_HSTS & false & false & true \\
SECURE\_COOKIES & false & false & true \\
TRUST\_PROXY & false & false & true \\
RATE\_LIMIT\_MODE & test & test & prod \\
JOB\_STORE\_BACKEND & memory & memory & redis \\
DEMO\_MODE & true & true & false \\
Step Timeout & 1200s (CPU) & 60s & Device-specific \\
Run Timeout & 1800s (CPU) & 120s & Device-specific \\
\hline
\end{tabular}
\end{table}

\subsection{Hardware-Specific Profiles}

Beyond environment modes, the platform supports hardware-specific deployment profiles through Docker Compose override files:

\paragraph{CPU Profile (.env.cpu)}
Optimized for CPU-only deployments with extended timeouts:

\begin{lstlisting}[language=bash, caption=CPU profile configuration]
# .env.cpu - CPU-optimized settings
LLM_DEVICE=cpu
LLM_STEP_TIMEOUT_SECONDS=1200
LLM_RUN_TIMEOUT_SECONDS=1800
LLM_MAX_CONCURRENT_CALLS=1
MEMGRAPH_BUILDER_LLM_TIMEOUT_MS=180000  # 3 minutes
\end{lstlisting}

\paragraph{GPU Profile (.env.gpu)}
Optimized for GPU-accelerated deployments:

\begin{lstlisting}[language=bash, caption=GPU profile configuration]
# .env.gpu - GPU-optimized settings
LLM_DEVICE=cuda
LLM_STEP_TIMEOUT_SECONDS=30
LLM_RUN_TIMEOUT_SECONDS=120
LLM_MAX_CONCURRENT_CALLS=4
MEMGRAPH_BUILDER_LLM_TIMEOUT_MS=10000  # 10 seconds
\end{lstlisting}

The Makefile provides convenient targets for launching with specific hardware profiles:

\begin{lstlisting}[language=bash, caption=Makefile targets for hardware profiles]
up-cpu: ## Start services with CPU profile
    $(DC) --env-file .env.cpu up -d --build

up-gpu: ## Start services with GPU profile
    $(DC) --env-file .env.gpu \
        -f docker-compose.yml \
        -f docker-compose.gpu.yml up -d --build
\end{lstlisting}

The GPU profile additionally includes Docker Compose GPU resource reservations:

\begin{lstlisting}[language=yaml, caption=GPU resource reservation in docker-compose.gpu.yml]
services:
  app:
    deploy:
      resources:
        reservations:
          devices:
            - driver: nvidia
              count: 1
              capabilities: [gpu]
    environment:
      NVIDIA_VISIBLE_DEVICES: all
      NVIDIA_DRIVER_CAPABILITIES: compute,utility
\end{lstlisting}

\subsection{Configuration Validation}

The platform validates configuration at startup to catch misconfiguration early:

\begin{enumerate}
    \item \textbf{Pydantic Validation}: All configuration values are validated against their type definitions
    \item \textbf{Field Validators}: Custom validators normalize and verify complex values
    \item \textbf{Computed Properties}: Derived values are calculated consistently
    \item \textbf{Health Probes}: The startup probe verifies database connectivity and dependency availability
\end{enumerate}

This multi-layered validation approach ensures that configuration errors are detected before the application accepts traffic, reducing the risk of runtime failures due to misconfiguration.

